This calculation concerns the precise periodic profile in Alp\"oge--Buckmaster,
\emph{Blowup for the Boussinesq equations with smooth forcing}, released
76-page source, Sections 3.1 and 4.1, equations (3.1), (3.2), (4.1)--(4.4).
It proves exact identities and a profile evolution for a prescribed affine
background, including the exact force that keeps the original profile.
It also identifies the change to the affine region when that force is
omitted. No conclusion about an infinite viscous controlled sequence follows
from these profile identities alone.

\section{The original profile and all coefficients}

Let $0<s_0<1$ and let $\chi\in C_c^\infty((-\pi,\pi))$ be even and equal to
one on $[-s_0,s_0]$. The actual source profile is the $2\pi$-periodic extension
of
\begin{equation}\label{pd:eq:F}
 F(s)=\sin s+\chi(s)(s-\sin s),\qquad -\pi<s<\pi.
\end{equation}
It is smooth because it equals $\sin s$ near the endpoints of this period;
it is odd because $\chi$ is even. Hence its period mean is zero. It equals
$s$ on $[-s_0,s_0]$, in particular $F'(0)=1$. Let $P$ be its unique
mean-zero periodic primitive. Indeed
$\widetilde P(s)=\int_0^s F(r)\,\PDd r$ is even and periodic, and
$P=\widetilde P-\PDmean{\widetilde P}$ has the claimed properties. Here
$\PDmean{h}=(2\pi)^{-1}\int_{-\pi}^{\pi}h(s)\,\PDd s$.

Fix a closed finite time interval $I=[t_0,t_1]$, smooth
$D:I\to\PDReal^{2\times2}$ with $\operatorname{tr}D=0$, smooth
$G:I\to\PDReal^2$, a frequency $\lambda>0$, and a nonzero initial vector
$\zeta(t_0)$. Define $J(x_1,x_2)=(-x_2,x_1)$ and
\begin{equation}\label{pd:eq:phase}
 \dot\zeta=-D^T\zeta,\qquad s(x,t)=\lambda\zeta(t)\cdot x,
 \qquad \kappa(t)=|\zeta(t)|^2,\qquad K(t)=\lambda^2\kappa(t).
\end{equation}
The notation $\kappa=|\zeta|^2$ is the source's geometric notation.
The physical temperature diffusivity is $\PDthermal\geq0$, and the physical
momentum viscosity is $\PDviscosity\geq0$; neither is the paper's time-scaling
parameter named $\nu$. These constants are not set to one. This calculation
does not rescale $A_0,\lambda_0,\mu$, the physical coordinates, time, or
laboratory gravity $e_2=(0,1)$.

The nonvanishing of $\zeta$ follows directly from the invertible fundamental
matrix of $\dot M=DM$, $M(t_0)=\mathrm{Id}$: differentiation gives
$\zeta=M^{-T}\zeta(t_0)$, and differentiation of the determinant gives
$\det M=\exp\int_{t_0}^t\operatorname{tr}D=1$.
Write
\begin{equation}\label{pd:eq:ab}
 a(t)=-\frac{J\zeta\cdot G}{\lambda\kappa},\qquad
 b(t)=\lambda\zeta_1.
\end{equation}
In particular the buoyancy coefficient uses the first laboratory component
$\zeta_1$.

For any smooth scalar $\theta$ and divergence-free velocity $u$, with
$\omega=\partial_1u_2-\partial_2u_1$, define the full retained-diffusion
residuals by
\begin{equation}\label{pd:eq:resdef}
 \mathcal R_\theta(\theta,u)
  =\partial_t\theta+u\cdot\nabla\theta-\PDthermal\Delta\theta,
 \qquad
 \mathcal R_\omega(\theta,u)
  =\partial_t\omega+u\cdot\nabla\omega-\partial_1\theta-\PDviscosity\Delta\omega.
\end{equation}
Thus these residuals are the required scalar force and curl of the required
momentum force. On the region under consideration the older fields are
$\theta_{\rm old}=G\cdot x$, $u_{\rm old}=Dx$; they need not have zero
residual. All statements below subtract their actual residuals.

\section{Exact residual and failure of constant-profile damping}

\begin{pdproposition}\label{pd:prop:res}
For arbitrary smooth real amplitudes $\Theta,\Omega$, define the source wave
\begin{equation}\label{pd:eq:wave}
 \vartheta=\Theta F(s),\quad
 \psi=\frac{\Omega}{\lambda^2\kappa}P(s),\quad
 v=\nabla^\perp\psi=\frac{\Omega}{\lambda\kappa}J\zeta F(s),\quad
 \varpi=\Delta\psi=\Omega F'(s).
\end{equation}
The scalar and curl residual increments are exactly
\begin{align}
 \mathcal R_\theta(\theta_{\rm old}+\vartheta,u_{\rm old}+v)
  -\mathcal R_\theta(\theta_{\rm old},u_{\rm old})
   &=(\dot\Theta-a\Omega)F-\PDthermal K\Theta F'',\label{pd:eq:rtheta}\\
 \mathcal R_\omega(\theta_{\rm old}+\vartheta,u_{\rm old}+v)
  -\mathcal R_\omega(\theta_{\rm old},u_{\rm old})
   &=(\dot\Omega-b\Theta)F'-\PDviscosity K\Omega F'''.\label{pd:eq:romega}
\end{align}
Every profile in these formulas is evaluated at $s(x,t)$.
\end{pdproposition}
\begin{proof}
Equation \eqref{pd:eq:phase} gives
$(\partial_t+Dx\cdot\nabla)s=
\lambda(\dot\zeta+D^T\zeta)\cdot x=0$.
The orthogonality $J\zeta\cdot\zeta=0$ gives both
$v\cdot\nabla\vartheta=0$ and $v\cdot\nabla\varpi=0$.
The older vorticity $D_{21}-D_{12}$ is spatially constant, so
$v\cdot\nabla\omega_{\rm old}=0$.
Moreover $v\cdot G=-a\Omega F$,
$\partial_1\vartheta=b\Theta F'$,
$\Delta\vartheta=K\Theta F''$, and
$\Delta\varpi=K\Omega F'''$.
Expanding \eqref{pd:eq:resdef}, canceling the older residual, and retaining all
these terms proves \eqref{pd:eq:rtheta}--\eqref{pd:eq:romega}.
The displayed velocity and vorticity in \eqref{pd:eq:wave} follow by
differentiating $P'=F$ twice; their signs follow from the specified $J$.
\end{proof}

The profile derivatives, including their transition-region terms, are
\begin{align}
 F'&=\cos s+\chi'(s)(s-\sin s)+\chi(s)(1-\cos s),\label{pd:eq:dF}\\
 F''&=(\chi-1)\sin s+2\chi'(1-\cos s)+\chi''(s-\sin s),\label{pd:eq:ddF}\\
 F'''&=(\chi-1)\cos s+3\chi'\sin s
              +3\chi''(1-\cos s)+\chi'''(s-\sin s).\label{pd:eq:dddF}
\end{align}
These follow by the product rule on $(-\pi,\pi)$ and extend smoothly
across the period endpoints.

\begin{pdproposition}\label{pd:prop:noeig}
There is no constant $c\in\PDReal$ such that $F''=cF$ on the circle, and no
constant $c\in\PDReal$ such that $F'''=cF'$ on the circle. Consequently, at
a time with $\PDthermal>0$ and $\Theta\ne0$, no choice of scalar
$\dot\Theta-a\Omega$ makes \eqref{pd:eq:rtheta} vanish for every phase.
At a time with $\PDviscosity>0$ and $\Omega\ne0$, no choice of scalar
$\dot\Omega-b\Theta$ makes \eqref{pd:eq:romega} vanish for every phase.
\end{pdproposition}
\begin{proof}
On $(-s_0,s_0)$, $F(s)=s$ and $F''=0$. An identity $F''=cF$ therefore
implies $cs=0$ for every such $s$, so $c=0$. Then $F''=0$ everywhere
would make $F$ globally affine. Its periodicity would force its slope to be
zero, contradicting $F'(0)=1$.
If $F'''=cF'$, evaluation inside that same interval gives $0=c$, since
$F'=1$ there. Then $F'''=0$ everywhere makes $F$ a quadratic polynomial
on $\PDReal$. Periodicity forces both its quadratic and linear coefficients
to vanish, again contradicting $F'(0)=1$.
Since $K>0$, each asserted residual cancellation with a nonzero amplitude
would imply precisely one of these excluded constant-coefficient
identities, by division by $\PDthermal K\Theta$ or $\PDviscosity K\Omega$.
\end{proof}

For example a proposed common damping equation
$\dot\Theta=a\Omega-d\Theta$, $\dot\Omega=b\Theta-d\Omega$ gives
$-d\Theta F-\PDthermal K\Theta F''$ and
$-d\Omega F'-\PDviscosity K\Omega F'''$ as its two residuals. On the exact affine
core they are $-d\Theta s$ and $-d\Omega$.
This proves failure of that fixed-profile amplitude prescription, not
failure of a viscous controlled construction.

\section{An exact evolving-profile map for all diffusivities}

We now construct the profile evolution instead of replacing the missing
profile dynamics by scalar damping. Let $T(s,t),V(s,t)$ be smooth,
mean-zero periodic functions, and let $Q_s=V$, $\PDmean Q=0$. Set
\begin{equation}\label{pd:eq:generalprofile}
 \vartheta=T(s,t),\quad \psi=\frac{Q(s,t)}{\lambda^2\kappa},\quad
 v=\frac{J\zeta}{\lambda\kappa}V(s,t),\quad \varpi=V_s(s,t).
\end{equation}
The time derivative $T_t$ here and below holds phase $s$ fixed.

\begin{pdproposition}\label{pd:prop:profile}
The residual increments for \eqref{pd:eq:generalprofile} are
\begin{equation}\label{pd:eq:profileR}
 T_t-aV-\PDthermal K T_{ss},\qquad
 \partial_s\big(V_t-bT-\PDviscosity K V_{ss}\big).
\end{equation}
Their vanishing is equivalent in the mean-zero class to
\begin{equation}\label{pd:eq:profilePDE}
 T_t=aV+\PDthermal K T_{ss},\qquad
 V_t=bT+\PDviscosity K V_{ss}.
\end{equation}
For the initial data $T(s,t_0)=\Theta_0F(s)$,
$V(s,t_0)=\Omega_0F(s)$, this system has a unique smooth periodic solution
on $I$ given by the following exact Fourier evolution. Write
\begin{equation}\label{pd:eq:fourier}
 F(s)=\sum_{n\geq1} f_n\sin(ns),\qquad
 f_n=\frac1\pi\int_{-\pi}^{\pi}F(s)\sin(ns)\,\PDd s.
\end{equation}
For each $n\geq1$, let $U_n(t,t_0)$ solve
\begin{equation}\label{pd:eq:Un}
 \partial_t U_n=
 \begin{pmatrix}-\PDthermal K n^2&a\\b&-\PDviscosity K n^2\end{pmatrix}U_n,
 \qquad U_n(t_0,t_0)=\mathrm{Id}.
\end{equation}
Then
\begin{equation}\label{pd:eq:profileSolution}
 \begin{pmatrix}T(s,t)\\V(s,t)\end{pmatrix}
 =\sum_{n\geq1}f_n U_n(t,t_0)
       \begin{pmatrix}\Theta_0\\\Omega_0\end{pmatrix}\sin(ns).
\end{equation}
\end{pdproposition}
\begin{proof}
The phase and tangency computations in Proposition~\ref{pd:prop:res} remain
valid, including $v\cdot\nabla T=v\cdot\nabla V_s=0$.
Now $v\cdot G=-aV$ and $\partial_1T=bT_s$, which proves
\eqref{pd:eq:profileR}. The expression inside its second derivative has zero
period mean. If its phase derivative vanishes, it is constant in $s$,
and its zero mean forces this constant to be zero. This proves the exact
equivalence to \eqref{pd:eq:profilePDE}.

For completeness, each matrix solution in \eqref{pd:eq:Un} can be defined by
the convergent iterated-integral series
$U_n=\mathrm{Id}+\sum_{j\geq1}\int_{t_0<r_j<\cdots<r_1<t}
 A_n(r_1)\cdots A_n(r_j)\,\PDd r_j\cdots\PDd r_1$, where $A_n$ is the
displayed matrix. For a fixed $n$, its $j$th term is bounded by
$(|I|\sup_I\|A_n\|)^j/j!$, so termwise differentiation gives the ODE.
Uniqueness follows by integrating the difference and applying the
iterated integral estimate, which tends to zero as $j\to\infty$.
For its column solution $y=(y_1,y_2)$, direct differentiation gives
\[
 \tfrac12\partial_t|y|^2
 =-Kn^2(\PDthermal y_1^2+\PDviscosity y_2^2)+(a+b)y_1y_2
 \leq \tfrac12(|a|+|b|)|y|^2.
\]
Multiplication by $\exp(-\int_{t_0}^t(|a|+|b|))$ and differentiation
therefore proves the bound, uniform in $n$,
\[
 |y(t)|\leq e^{\frac12\int_{t_0}^t(|a|+|b|)}|y(t_0)|.
\]
Repeated differentiation of \eqref{pd:eq:Un} gives, for each fixed $j$,
$\sup_I|\partial_t^j y|\leq C_j(1+n)^{2j}|y(t_0)|$, by induction:
the coefficient derivatives of $A_n$ are bounded by constants times
$(1+n)^2$, and the differentiated product has finitely many terms.

Integration by parts any prescribed number of times in \eqref{pd:eq:fourier}
shows $|f_n|\leq C_N n^{-N}$ for every $N$, without endpoint terms by
periodicity. Thus \eqref{pd:eq:profileSolution} and every mixed derivative
converge uniformly on $\PDTorus\times I$, and termwise differentiation proves
the equations and the prescribed initial values. The Fourier expansion
indeed equals $F$: its absolutely convergent series has the same Fourier
coefficients as $F$, and a continuous function with all Fourier
coefficients zero is zero. To see this last assertion directly, convolve
with the Fej\'er kernel
$\mathfrak F_N(r)=N^{-1}(\sin(Nr/2)/\sin(r/2))^2$.
It has integral $2\pi$, is nonnegative, and its integral outside every
fixed neighborhood of $0$ tends to zero because its numerator is at
most one and the denominator is bounded below there. Uniform continuity
therefore gives uniform convergence of its convolution to the continuous
function. When all Fourier coefficients are zero each such convolution
is zero, proving the assertion. The same Fourier coefficient ODE and
its uniqueness show uniqueness among smooth solutions of
\eqref{pd:eq:profilePDE}.
\end{proof}

\begin{pdproposition}\label{pd:prop:pressure}
The map in Proposition~\ref{pd:prop:profile} also realizes the full momentum
equation. Let $S_s=T$, $\PDmean S=0$, and retain $Q_s=V$, $\PDmean Q=0$.
For a solution of \eqref{pd:eq:profilePDE}, the exact pressure increment is
\begin{equation}\label{pd:eq:pressure}
 p^{\rm wave}(x,t)=
 -\frac{2\zeta\cdot D J\zeta}{\lambda^2\kappa^2}Q(s,t)
 +\frac{\zeta_2}{\lambda\kappa}S(s,t).
\end{equation}
With this pressure increment, adding $(T,v)$ to any affine older state
preserves its actual vector momentum residual, including retained
viscosity $\PDviscosity$ and the unchanged laboratory buoyancy $\theta e_2$.
\end{pdproposition}
\begin{proof}
Put $c=J\zeta/(\lambda\kappa)$, a vector in this proof. The vector
momentum residual increment before adding pressure is
\[
 (\dot c+Dc)V+c(V_t-\PDviscosity K V_{ss})-Te_2.
\]
Indeed the transported phase has zero material derivative under $Dx$,
$v\cdot\nabla v=0$, $v\cdot\nabla(Dx)=Dv$, and
$\Delta v=KcV_{ss}$.
For every trace-free $2\times2$ matrix,
$DJ=-JD^T$; multiplying its four entries verifies the identity.
It gives $\dot c=Dc-c\dot\kappa/\kappa$ and
$\dot\kappa=-2\zeta\cdot D\zeta$.
The identity
$(J\zeta)\cdot D(J\zeta)+\zeta\cdot D\zeta
=\kappa\operatorname{tr}D=0$ shows that $\dot c+Dc$ is orthogonal to
$J\zeta$. Its component along $\zeta$ is therefore
\[
 \dot c+Dc=\frac{2\zeta\cdot D J\zeta}{\lambda\kappa^2}\zeta.
\]
Also $bc-e_2=-\zeta_2\zeta/\kappa$, by the two explicit components of
$J\zeta=(-\zeta_2,\zeta_1)$. Substituting the second profile equation
reduces the residual to
\[
 \zeta\left[
 \frac{2\zeta\cdot D J\zeta}{\lambda\kappa^2}V
 -\frac{\zeta_2}{\kappa}T\right].
\]
The spatial gradient of \eqref{pd:eq:pressure} is its negative, because
$\nabla s=\lambda\zeta$, $Q_s=V$ and $S_s=T$.
This proves the full momentum statement. Mean-zero periodic primitives
exist because the profiles have zero mean, by the same integration
argument used for $P$ above.
\end{proof}

\begin{pdproposition}\label{pd:prop:heat}
If $\PDthermal=\PDviscosity=\delta\geq0$, put
\begin{equation}\label{pd:eq:heattime}
 \tau(t)=\delta\lambda^2\int_{t_0}^t|\zeta(r)|^2\,\PDd r,
 \qquad
 F_{\tau}(s)=\PDheat_{\tau}F(s)=\sum_{n\geq1}f_ne^{-n^2\tau}\sin(ns).
\end{equation}
Let $\dot\Theta=a\Omega$, $\dot\Omega=b\Theta$ with the given initial
amplitudes. The exact solution in Proposition~\ref{pd:prop:profile} is
\begin{equation}\label{pd:eq:heatmap}
 T=\Theta F_\tau,\quad V=\Omega F_\tau,\quad
 Q=\Omega\PDheat_\tau P,
 \quad \varpi=\Omega\partial_sF_\tau.
\end{equation}
In particular this is an exact retained-diffusion map with the same
prescribed $D,G,\zeta$, all original frequency factors, and laboratory
gravity. Each nonzero temperature mode has exact gain
\begin{equation}\label{pd:eq:gain}
 \frac{T_n(t)}{T_n(t_0)}
   =\frac{\Theta(t)}{\Theta_0}e^{-n^2\tau(t)}
 \quad(f_n\Theta_0\ne0).
\end{equation}
For vorticity the corresponding ratio is
$(\Omega(t)/\Omega_0)e^{-n^2\tau(t)}$ when $nf_n\Omega_0\ne0$.
\end{pdproposition}
\begin{proof}
The matrix in \eqref{pd:eq:Un} is $B(t)-\delta K(t)n^2\mathrm{Id}$,
where $B=\left(\begin{smallmatrix}0&a\\b&0\end{smallmatrix}\right)$.
If $U$ solves $\dot U=BU$, differentiation of
$e^{-n^2\tau(t)}U(t,t_0)$ gives exactly \eqref{pd:eq:Un}, including its
initial identity. Hence uniqueness proves this factorization and
\eqref{pd:eq:heatmap}. Since derivatives commute termwise with the Fourier
series, $(\PDheat_\tau P)_s=\PDheat_\tau F$ and its mean remains zero.
This verifies the streamfunction as well as the scalar and vorticity
profiles. Formula \eqref{pd:eq:gain} is division of the actual mode
coefficients, retaining its stated nonzero denominator.
\end{proof}

The exact temperature phase norm is
\begin{equation}\label{pd:eq:norm}
 \|T(\cdot,t)\|_{L^2(-\pi,\pi)}^2
 =\pi\Theta(t)^2\sum_{n\geq1}f_n^2e^{-2n^2\tau(t)}.
\end{equation}
This follows by integrating the uniformly convergent squared series and
using the exact orthogonality identity
\[
 \int_{-\pi}^{\pi}\sin(ns)\sin(ms)\,\PDd s=\pi\mathbf1_{n=m}.
\]
Thus $\|F_\tau\|_2\leq e^{-\tau}\|F\|_2$, and for any index with
$f_n\ne0$ one has $\|F_\tau\|_2\geq\sqrt\pi|f_n|e^{-n^2\tau}$.
These are inequalities about phase profiles; an unlocalized affine wave
need not be integrable over $\PDReal^2$.

\section{Positive heat time and the affine core}

\begin{pdproposition}\label{pd:prop:analytic}
For every $\tau>0$, $\PDheat_\tau F$ has no affine restriction to any nonempty
open interval of $\PDReal$. In particular it cannot retain the source's
nontrivial exact linear phase interval. Also $F$ has infinitely many
nonzero Fourier modes.
\end{pdproposition}
\begin{proof}
For any $R>0$ and any nonnegative integer $m$, the series of $m$th complex
derivatives of
$\sum_{n\geq1}f_ne^{-n^2\tau}\sin(nz)$ is bounded on
$|\operatorname{Im}z|\leq R$ by
$\sum_{n\geq1}|f_n|n^me^{-n^2\tau+nR}<\infty$.
For instance boundedness of $f_n$ and completion of the square make the
tail bounded by a polynomial times $e^{-\tau n^2/2}$.
The sum therefore defines an entire function $h(z)$ and has its
termwise derivatives. If $h(s)=cs+d$ on a nonempty real open interval,
the entire function $h(z)-cz-d$ vanishes identically. The identity
principle used here follows from Taylor expansion: a nonzero holomorphic
function has a first nonzero Taylor coefficient at each zero where its
Taylor series is not identically zero, hence that zero is isolated;
an accumulating set of zeros forces the Taylor series to vanish locally,
and this local vanishing propagates through the connected domain.
Periodicity now gives $2\pi c=h(z+2\pi)-h(z)=0$. Oddness gives $d=0$.
Thus every Fourier coefficient $f_ne^{-n^2\tau}$ would vanish, implying
$f_n=0$ for all $n$. The Fourier uniqueness argument in
Proposition~\ref{pd:prop:profile} would imply $F=0$, contradicting $F'(0)=1$.
If $F$ had only finitely many nonzero modes, the same entire-function
argument, applied to its finite trigonometric sum at heat time zero,
would contradict its interval $F(s)=s$.
\end{proof}

There is nevertheless an exact estimate of the defect deep inside the
original interval. Define, for $\tau>0$ and $d>0$,
\begin{equation}\label{pd:eq:tail}
 \mathcal E(d,\tau)=\frac{2}{\sqrt\pi}
          \int_{d/(2\sqrt\tau)}^\infty e^{-r^2}\,\PDd r
 \leq\frac{2\sqrt\tau}{d\sqrt\pi}e^{-d^2/(4\tau)}.
\end{equation}
The inequality follows from $r\geq d/(2\sqrt\tau)$ and
$\int_a^\infty r e^{-r^2}\PDd r=e^{-a^2}/2$.

\begin{pdproposition}\label{pd:prop:tail}
For $0<r_0<s_0$, $d=s_0-r_0$, and $|s|\leq r_0$,
\begin{align}
 |\partial_sF_\tau(s)-1|&\leq\|F'-1\|_\infty\mathcal E(d,\tau),\label{pd:eq:tail1}\\
 |F_\tau(s)-s|&\leq |s|\|F'-1\|_\infty\mathcal E(d,\tau),\label{pd:eq:tail0}\\
 |\partial_s^mF_\tau(s)|&\leq\|F^{(m)}\|_\infty\mathcal E(d,\tau)
   \quad(m\geq2).\label{pd:eq:tailm}
\end{align}
\end{pdproposition}
\begin{proof}
The periodized Gaussian kernel of $\PDheat_\tau$ is
$\sum_{k\in\mathbb Z}(4\pi\tau)^{-1/2}
e^{-(r+2\pi k)^2/(4\tau)}$. To verify this representation, its $n$th
Fourier multiplier is the Gaussian integral
$\int_{\PDReal}(4\pi\tau)^{-1/2}e^{-r^2/(4\tau)}e^{-inr}\PDd r=e^{-n^2\tau}$.
One direct verification of this integral differentiates it with respect
to $n$, integrates the Gaussian derivative by parts, and obtains
$I'(n)=-2\tau n I(n)$ with $I(0)=1$; the last normalization is the usual
Gaussian integral, obtained by squaring it and integrating in polar
coordinates. Fourier uniqueness then proves the representation.
For any periodic bounded function $h$ vanishing on $[-s_0,s_0]$, its
periodic lift satisfies $h(s-r)=0$ whenever $|s|\leq r_0$ and $|r|<d$.
The Gaussian representation on the line therefore gives
\[
 |\PDheat_\tau h(s)|\leq
 \|h\|_\infty\int_{|r|\geq d}
 (4\pi\tau)^{-1/2}e^{-r^2/(4\tau)}\,\PDd r
 =\|h\|_\infty\mathcal E(d,\tau).
\]
Apply this to $h=F'-1$ and $h=F^{(m)}$ for $m\geq2$, using commutation
of derivatives with $\PDheat_\tau$ and $\PDheat_\tau1=1$.
Finally $F_\tau(0)=0$ by oddness; integrating the first bound between
$0$ and $s$ proves the second bound with the same $d$.
\end{proof}

The source background in equation (3.3) is obtained because each older
layer is exactly affine on the newer support, by equations (4.3)--(4.6).
For a heat-evolved older layer with nonzero scalar or velocity amplitude,
Proposition~\ref{pd:prop:analytic} removes this exact affine premise.
Its gradient at the origin can still be computed, but its values away
from the origin are not that gradient times $x$. The tail estimates give
explicit defects rather than an identification of those two fields.
Consequently Proposition~\ref{pd:prop:heat} alone supplies no coupled feedback
law for the heat-modified older background.

\section{Exact preservation by force and the localized diffusion costs}

For the original inviscid amplitude equations $\dot\Theta=a\Omega$,
$\dot\Omega=b\Theta$, Proposition~\ref{pd:prop:res} shows that the exact
additional scalar and curl forces are
\begin{equation}\label{pd:eq:repair}
 f_{\theta,\mathrm{diff}}=-\PDthermal K\Theta F'',\qquad
 E_{\mathrm{diff}}=-\PDviscosity K\Omega F'''.
\end{equation}
They vanish on the entire affine interval $|s|<s_0$. This retains exactly
the original core, amplitudes, phase and hence its local affine feedback;
it pays a nonzero diffusion force in the rest of the profile. In the
unlocalized wave this force is periodic in phase and is not a compactly
supported force on $\PDReal^2$.

Here is the exact compact counterpart for the leading fields in source
equation (4.2). Let $M$ solve $\dot M=DM$, $M(t_0)=\mathrm{Id}$, let
$p=\zeta(t_0)$, and let $\ell>0$. Fix the source's radial
$f\in C_c^\infty(B_1)$, $f=1$ on $B_{1/2}$, and put
\begin{align}\label{pd:eq:localized}
 a_{\rm mat}&=M^{-1}x,\qquad g(a_{\rm mat})=f(a_{\rm mat}/\ell),\notag\\
 \theta^L&=w\Theta F(s)g(a_{\rm mat}),\qquad
 \psi^L=\frac{w\Omega}{\lambda^2\kappa}P(s)g(a_{\rm mat}),\quad
 u^L=\nabla^\perp\psi^L.
\end{align}
The activation $w(t)$ remains the original smooth scalar function; no
time differentiation of it is omitted from the existing inviscid
residual. The following formulas give exactly the additional diffusion
terms, with $w$ retained.
At a fixed time put
\begin{equation}\label{pd:eq:L}
 B=M^{-1},\quad H=BB^T,\quad q=B\zeta,\quad
 \mathscr L=\lambda^2\kappa\partial_s^2+
            2\lambda q\cdot\nabla_a\partial_s+H:\nabla_a^2.
\end{equation}
Although $s=\lambda p\cdot a_{\rm mat}$ upon evaluation, $s,a$ are
independent variables in this differential operator.

\begin{pdproposition}\label{pd:prop:localized}
For every smooth profile $A(s,a)$,
$\Delta_x[A(\lambda\zeta\cdot x,Bx)]=
(\mathscr L A)(\lambda\zeta\cdot x,Bx)$.
Consequently the exact added forces for \eqref{pd:eq:localized} are
\begin{align}
 f^L_{\theta,\mathrm{diff}}
 &=-\PDthermal w\Theta\big[
 \lambda^2\kappa F''g+2\lambda F' q\cdot\nabla g
                         +F(H:\nabla^2g)\big],\label{pd:eq:localtheta}\\
 f^L_{u,\mathrm{diff}}&=-\PDviscosity\Delta u^L
 =-\PDviscosity\nabla^\perp\Delta\psi^L,
 \qquad
 E^L_{\mathrm{diff}}=-\PDviscosity\frac{w\Omega}{\lambda^2\kappa}
       \mathscr L^2(Pg).\label{pd:eq:localomega}
\end{align}
Writing $\mathscr H=H:\nabla_a^2$, the last operator is explicitly
\begin{align}\label{pd:eq:L2}
 \mathscr L^2(Pg)={}&
 \lambda^4\kappa^2 F'''g
 +4\lambda^3\kappa F''q\cdot\nabla g
 +2\lambda^2\kappa F'\mathscr H g
 +4\lambda^2F'(q\cdot\nabla)^2g\notag\\
 &+4\lambda F(q\cdot\nabla)\mathscr H g
 +P\mathscr H^2g.
\end{align}
All these forces are supported in $M\overline B_\ell$. They vanish in
the exact affine region $M B_{\ell/2}\cap\{|s|<s_0\}$.
\end{pdproposition}
\begin{proof}
The chain rule gives
$\partial_{x_i}A=\lambda\zeta_i A_s+
\sum_jB_{ji}A_{a_j}$. Applying it a second time and summing in $i$
gives the three terms in \eqref{pd:eq:L}, since
$\sum_i B_{ji}\zeta_i=(B\zeta)_j$ and
$\sum_iB_{ji}B_{ki}=(BB^T)_{jk}$. Applying the same identity twice to
$Pg$ proves \eqref{pd:eq:localomega}. Expanding the square of the
commuting constant-in-$(s,a)$ operators in \eqref{pd:eq:L}, and using
$P'=F$, $P''=F'$, $P'''=F''$, $P''''=F'''$, proves all six terms
of \eqref{pd:eq:L2}.
Spatial differentiation does not enlarge the support of a smooth
compactly supported function, proving the support statement.
On the plateau $g=1$ all positive derivatives of $g$ vanish.
On the linear phase interval $F''=F'''=0$. Thus the scalar and curl
formulas vanish on their intersection. There
$\psi^L=(w\Omega/(\lambda^2\kappa))(P(0)+s^2/2)$ is quadratic in $x$,
so $u^L$ is linear and $\Delta u^L=0$, proving the vector-force statement
as well. Finally $\operatorname{curl}(-\PDviscosity\Delta u^L)
=-\PDviscosity\Delta^2\psi^L$, so the specified vector force realizes the
curl force without any inverse-curl ambiguity.
\end{proof}

For explicit quantitative costs let
$C_j=\sup_a\|\nabla^jf(a)\|_{\rm HS}$, with the Euclidean
Hilbert--Schmidt tensor norm, and $C_0=\|f\|_\infty$.
Then $\|\nabla^jg\|_\infty=C_j\ell^{-j}$ by the chain rule. Write
$h=\|H\|_{\rm HS}$ and $q_*=|q|$. Cauchy--Schwarz on the tensor
contractions in \eqref{pd:eq:localtheta} and \eqref{pd:eq:L2} proves
\begin{align}
 \|f^L_{\theta,\mathrm{diff}}\|_\infty
 \leq{}&\PDthermal|w\Theta|\big[
 \lambda^2\kappa\|F''\|_\infty C_0
 +2\lambda q_*\|F'\|_\infty C_1\ell^{-1}
 +h\|F\|_\infty C_2\ell^{-2}\big],\label{pd:eq:costtheta}\\
 \|E^L_{\mathrm{diff}}\|_\infty
 \leq{}&\frac{\PDviscosity|w\Omega|}{\lambda^2\kappa}\big[
 \lambda^4\kappa^2\|F'''\|_\infty C_0
 +4\lambda^3\kappa q_*\|F''\|_\infty C_1\ell^{-1}\notag\\
 &\quad+(2\lambda^2\kappa h+4\lambda^2q_*^2)
       \|F'\|_\infty C_2\ell^{-2}
 +4\lambda q_*h\|F\|_\infty C_3\ell^{-3}
 +h^2\|P\|_\infty C_4\ell^{-4}\big].\label{pd:eq:costomega}
\end{align}
The same reasoning applies to every compact correction: for a specified
smooth correction $(\theta^c,\psi^c)$ the extra scalar and vector forces
are exactly $-\PDthermal\Delta\theta^c$ and
$-\PDviscosity\nabla^\perp\Delta\psi^c$, with curl
$-\PDviscosity\Delta^2\psi^c$. This is the exact linear map from an already
constructed inviscid forced field to a retained-diffusion forced field:
\begin{equation}\label{pd:eq:force-map}
 (\theta,u,p,f_\theta,f_u)\longmapsto
 (\theta,u,p,f_\theta-\PDthermal\Delta\theta,
                      f_u-\PDviscosity\Delta u).
\end{equation}
Substitution in the scalar and momentum equations proves the map:
adding $\PDthermal\Delta\theta$ or $\PDviscosity\Delta u$ to the new right-hand side
cancels its respective force decrement. Thus pressure, divergence,
initial data and the entire trajectory are retained exactly. On each
closed finite smooth stage these are smooth finite forces; where the
original fields and their derivatives are compactly supported, the new
terms have the same support. This statement supplies no uniform bound
as the layer index tends to infinity. The explicit factors
$\lambda^2$, $\ell^{-2}$, and $\ell^{-4}$ above must be retained in such
an analysis.

\section*{Source and scope receipt}
The primary source read here is the 76-page PDF
\emph{Blowup for the Boussinesq equations with smooth forcing},
Levent Alp\"oge and Tristan Buckmaster; source equations (3.1)--(3.6)
are on pages 8--9 and (4.1)--(4.6) on pages 37--38.
The exact file SHA-256 is
\begin{center}\small\ttfamily
895a628d1783bcb039374686f50b895b5f450f53b8ef8aa173523487a7a4a21b.
\end{center}
The companion replay script checks the profile differentiation, full
two-dimensional residual algebra, localized Laplacian square, common
diffusivity matrix factorization, and eigenfunction obstruction
identities. Its checks supplement the proofs above; they do not certify
an infinite-stage existence or force-smoothness theorem.
